{\footnotesize
\begin{threeparttable}
\begin{tabular}{lrrrrrrrrr}
\toprule
 & \multicolumn{3}{c}{2014} & \multicolumn{3}{c}{2014 to 2019} & \multicolumn{3}{c}{2019 to 2022} \\
\cmidrule(lr){2-4}\cmidrule(lr){5-7}\cmidrule(lr){8-10}
State & Jobs & Wage (MW) & Portab. & $T$ & $D$ & $F$ & $T$ & $D$ & $F$ \\
\midrule
Rio de Janeiro & 3,749,182 & 3.18 & 0.677 & 0.071 & 0.024 & 2.94 & 0.050 & 0.020 & 2.46 \\
Rio Grande do Sul & 2,570,628 & 2.55 & 0.653 & 0.048 & 0.017 & 2.83 & 0.038 & 0.013 & 2.80 \\
São Paulo & 12,205,363 & 3.24 & 0.693 & 0.053 & 0.019 & 2.79 & 0.039 & 0.015 & 2.60 \\
Paraná & 2,628,186 & 2.46 & 0.679 & 0.056 & 0.020 & 2.75 & 0.044 & 0.031 & 1.40 \\
Minas Gerais & 4,049,204 & 2.24 & 0.659 & 0.049 & 0.019 & 2.57 & 0.040 & 0.016 & 2.51 \\
Bahia & 1,735,685 & 2.15 & 0.652 & 0.065 & 0.030 & 2.15 & 0.058 & 0.045 & 1.31 \\
Espírito Santo & 774,335 & 2.34 & 0.668 & 0.055 & 0.029 & 1.92 & 0.056 & 0.043 & 1.30 \\
Sergipe & 293,384 & 2.11 & 0.638 & 0.101 & 0.055 & 1.81 & 0.057 & 0.046 & 1.23 \\
Mato Grosso & 644,535 & 2.30 & 0.657 & 0.063 & 0.039 & 1.64 & 0.051 & 0.034 & 1.50 \\
Distrito Federal & 742,729 & 2.93 & 0.666 & 0.089 & 0.055 & 1.60 & 0.065 & 0.053 & 1.22 \\
Pernambuco & 1,330,510 & 2.05 & 0.671 & 0.088 & 0.056 & 1.59 & 0.047 & 0.024 & 1.97 \\
Goiás & 1,175,134 & 2.22 & 0.667 & 0.064 & 0.041 & 1.57 & 0.055 & 0.052 & 1.06 \\
Mato Grosso do Sul & 505,598 & 2.18 & 0.636 & 0.080 & 0.055 & 1.45 & 0.053 & 0.121 & 0.44 \\
Paraíba & 402,422 & 1.74 & 0.652 & 0.097 & 0.067 & 1.44 & 0.079 & 0.080 & 0.99 \\
Ceará & 1,135,961 & 1.84 & 0.653 & 0.082 & 0.059 & 1.40 & 0.067 & 0.120 & 0.56 \\
Santa Catarina & 1,977,166 & 2.43 & 0.684 & 0.062 & 0.048 & 1.29 & 0.036 & 0.014 & 2.60 \\
Pará & 750,506 & 2.27 & 0.660 & 0.097 & 0.078 & 1.25 & 0.051 & 0.053 & 0.96 \\
Rio Grande do Norte & 432,497 & 1.87 & 0.647 & 0.094 & 0.076 & 1.24 & 0.067 & 0.061 & 1.10 \\
Piauí & 275,278 & 1.76 & 0.651 & 0.092 & 0.077 & 1.20 & 0.059 & 0.082 & 0.72 \\
Rondônia & 245,342 & 2.12 & 0.660 & 0.122 & 0.112 & 1.09 & 0.051 & 0.035 & 1.44 \\
Alagoas & 349,474 & 1.76 & 0.648 & 0.110 & 0.127 & 0.87 & 0.059 & 0.058 & 1.03 \\
Maranhão & 462,281 & 2.02 & 0.658 & 0.099 & 0.121 & 0.82 & 0.077 & 0.375 & 0.20 \\
Amapá & 75,501 & 2.08 & 0.651 & 0.146 & 0.199 & 0.73 & 0.104 & 0.214 & 0.49 \\
Tocantins & 154,041 & 1.96 & 0.633 & 0.096 & 0.186 & 0.52 & 0.061 & 0.077 & 0.80 \\
Roraima & 46,007 & 1.92 & 0.657 & 0.093 & 0.321 & 0.29 & 0.090 & 0.177 & 0.51 \\
Acre & 75,589 & 1.94 & 0.647 & 0.094 & 0.667 & 0.14 & 0.098 & 1.579 & 0.06 \\
Amazonas & 437,852 & 2.46 & 0.697 & 0.094 & 2.325 & 0.04 & 0.051 & 0.058 & 0.88 \\
\bottomrule
\end{tabular}
\begin{tablenotes}[flushleft]\footnotesize
\item \textit{Note:} States ranked by allocative flexibility from 2014 to 2019. Source: RAIS, authors' calculations.
\end{tablenotes}
\end{threeparttable}
}
